% =====================================================================
\chapter{Introduction}\label{ch:intro}
% =====================================================================

% \begin{plan}[Chapter plan --- Introduction]
% \pages{5 pages.}\quad
% \textbf{Job of this chapter:} convince the reader there is a real measurement
% problem (audits read what a model \emph{says}), state the RQs exactly, and list
% contributions that match what Chapters 5--6 actually show. No equations
% except a pointer to \cref{ch:framework}. One figure (\cref{fig:where-measured}).
% \end{plan}

% ---------------------------------------------------------------------
\section{Background and motivation}\label{sec:intro-motivation}
% \begin{plan}
% \begin{itemize}
%   \item VLMs are now used as automated judges for retrieval and moderation
%         (``MLLM-as-a-judge'') \citeplan{chen2024mllmjudge,zheng2023judging}.
%   \item Image search shapes beliefs about the world: occupational
%         under-representation and its effect on users
%         \cite{kay2015unequal,metaxa2021image}; warm/competent trait stereotypes
%         in results \cite{otterbacher2017competent}.
%   \item VLMs inherit these stereotypes, and image features contribute more
%         than text \cite{jiang2024modscan,weng2024images,hall2023visogender};
%         intrinsic encoder bias propagates to zero-shot retrieval
%         \cite{ghate2025biases}.
%   \item Close with the human anchor: people who score higher on the ASI are
%         less likely to see gender bias in image results
%         \cite{otterbacher2018sigir} --- so ``does the judge see bias like a
%         sexist or a non-sexist human?'' is a measurable question.
% \end{itemize}
% \end{plan}

Vision language models (VLMs) are increasingly used to evaluate images in
retrieval and moderation systems. These models can produce quality and
relevance judgements at a scale that is difficult to achieve through human
annotation, motivating their use as multimodal judges
\cite{chen2024mllmjudge,zheng2023judging}. Their suitability for this role
depends on whether the judgements are reliable and whether they reproduce
social biases present in the training data.

Gender bias in image search provides a practical setting for examining these
questions. Search results can underrepresent women in occupations, and exposure
to those results can shift users' beliefs about occupational gender distributions
\cite{kay2015unequal}. Later audits document representational inequalities
across platforms and time \cite{metaxa2021image}. The associations also extend
to personal traits, with images of women linked more closely to warmth than
competence \cite{otterbacher2017competent}.

Related concerns arise when learned visual representations are used to judge
such content. Studies report stereotyped associations in multimodal models and
identify contributions from visual features under the evaluated conditions
\cite{jiang2024modscan,weng2024images}. Models can also resolve gendered
references in occupational images along stereotyped lines
\cite{hall2023visogender}, while associations measured in vision encoders can
carry into retrieval rankings \cite{ghate2025biases}. These findings motivate
evaluating automated judges instead of assuming that replacing human raters
removes bias from the evaluation.

A human study provides a reference for this comparison. Otterbacher et al.\
\cite{otterbacher2018sigir} administered the Ambivalent Sexism Inventory
(ASI) \cite{glick1996asi} and asked participants to evaluate the objectivity
of image search result grids. Participants higher in benevolent sexism were
less likely to identify particular gender-skewed grids as non-objective.
This thesis examines whether inventory responses and search judgements show
a corresponding relationship across VLMs. The comparison uses the historical
images and human ratings, although the models rate individual images rather
than complete grids, and the available baseline groups use total ASI.

% ---------------------------------------------------------------------
\section{Problem statement and measurement reliability}\label{sec:intro-problem}
% \begin{plan}
% \begin{itemize}
%   \item Alignment training changes what models say
%         \cite{ouyang2022training,bai2022constitutional}.
%   \item Evidence that it hides rather than removes bias: overt stereotypes
%         become positive while covert ones stay negative
%         \cite{hofmann2024covert}; explicitly unbiased models still form
%         biased associations \cite{bai2025explicitly}; measured gender bias
%         changes when the prompt signals an evaluation \cite{gao2025measuring};
%         persona framing changes social reasoning \cite{yang2026persona}.
%   \item Define \textbf{algorithmic defensiveness} (the thesis's term) in one
%         sentence: a sanitised response triggered by evaluative framing or, as
%         found here, by the mere presence of an image.
%   \item Using a model as the auditor adds its own position and scoring biases
%         \cite{zheng2023judging,wang2024fair,ye2025justice}.
%   \item The gap: no audit measures a multimodal model's distribution
%         \emph{before} sampling \emph{and} anchors it to a validated human scale.
% \end{itemize}
% \end{plan}

Direct questions about gender attitudes require care because alignment
training shapes the responses a model produces. Instruction tuning and
reinforcement learning from human feedback modify response behaviour
\cite{ouyang2022training}, while constitutional objectives train responses
against explicit principles \cite{bai2022constitutional}. Whether these
procedures also change the social associations measured by an audit is a
separate empirical question. Several studies distinguish explicit responses
from indirectly elicited associations.

Hofmann et al.\ \cite{hofmann2024covert} found that in models trained with
human feedback, overtly elicited stereotypes became markedly more favourable
while covertly elicited ones stayed negative, so alignment improved the self
report rather than the underlying association. Models that describe themselves
as unbiased still form associations consistent with stereotypes when probed
indirectly \cite{bai2025explicitly}. Gao and Kreiss \cite{gao2025measuring}
showed that measured gender bias in large language models is brittle with
respect to framing: measured bias scores differ depending on whether the
prompt reads as a test. Work on persona framing
reports the same sensitivity in social reasoning tasks \cite{yang2026persona}.
These findings motivate checking whether audit framing changes the responses
being measured before interpreting those responses as evidence of attitudes.

This thesis uses algorithmic defensiveness to describe the hypothesis that
evaluative cues can alter a model's response behaviour and obscure
gender-related associations. The experiments examine response changes across
prompt and image conditions, including the featureless grey patch
(\cref{sec:s1-quality}). These comparisons assess measurement sensitivity;
reduced response coverage alone does not identify refusal or deliberate
concealment.

Using a second model to evaluate responses introduces additional measurement
concerns. Judges built on large language models carry their own
position and verbosity biases together with a preference for their own outputs
\cite{zheng2023judging}, their scores are sensitive to the order and
presentation of the options they are given \cite{wang2024fair}, and these
effects are systematic and quantifiable rather than incidental
\cite{ye2025justice}. An audit built on a judge inherits the judge's artefacts.

The gap addressed by this thesis concerns the combination of probability
measurements, a human questionnaire and image search judgements. The audit
reads the distribution over the next token in a single forward pass and
compares the resulting measurements with the ASI and SIGIR human baselines.
This avoids sampling and text parsing; it does not remove the effects of
alignment training from the model weights or their output distribution.
\cref{fig:where-measured} shows
where in the pipeline each kind of audit takes its reading.

\begin{figure}[htbp]\centering
\resizebox{\linewidth}{!}{%
\begin{tikzpicture}[node distance=7mm and 9mm,
  box/.style={draw,rounded corners,align=center,font=\small,minimum height=9mm,inner sep=4pt},
  arr/.style={-{Stealth[length=2.5mm]},thick}]
\node[box] (in) {Prompt $x$\\+ image $v$};
\node[box,right=of in] (fwd) {Single\\forward pass};
\node[box,right=of fwd,fill=blue!8] (dist) {Next-token\\distribution $p(\cdot\mid x,v)$};
\node[box,right=of dist] (samp) {Sampling /\\decoding};
\node[box,right=of samp] (txt) {Generated\\text};
\draw[arr] (in) -- (fwd); \draw[arr] (fwd) -- (dist);
\draw[arr] (dist) -- (samp); \draw[arr] (samp) -- (txt);
\node[below=6mm of dist,font=\footnotesize,align=center,text=blue!60!black]
  (this) {\textbf{this thesis measures here}\\($\pyes$, $\capm$, $\rexp$)};
\node[below=6mm of txt,font=\footnotesize,align=center,text=red!60!black]
  (usual) {self-report audits\\read here};
\draw[-{Stealth},blue!60!black] (this) -- (dist);
\draw[-{Stealth},red!60!black] (usual) -- (txt);
\end{tikzpicture}}
\caption{Where the measurement is taken. The audit records next-token
probabilities before decoding and reports how much mass falls within the
intended response set. Other tokens may begin explanations, refusals or
alternative response formats, which captured mass alone cannot distinguish.%
\ifshowplan\ (Draft TikZ, restyle as needed.)\fi}
\label{fig:where-measured}
\end{figure}

% ---------------------------------------------------------------------
\section{Research questions and aims}\label{sec:intro-rq}
% \begin{plan}
% State the RQs verbatim (they match the proposal and the thesis plan):
% \begin{description}
%   \item[RQ1] What is the intrinsic psychological baseline attitude of the model
%         when evaluated in a vacuum?
%   \item[RQ2] Is that intrinsic distributional skew descriptively consistent
%         with the model's extrinsic, downstream behaviour when evaluating
%         real-world search queries?
% \end{description}
% Then the four proposal aims, each with one line saying how far it was met
% (Aim 1 met for 2 production models; Aim 2 met on the historical stream only;
% Aim 3 descriptive only, $n=2$; Aim 4 met via $\dRR$). Note that RQ2 is
% stated descriptively because two models cannot support a correlation.
% \end{plan}

The proposal framed the study in terms of an intrinsic psychological baseline
and its relationship with downstream behaviour. The delivered experiments
operationalise these aims as two questions about measured responses:

\begin{description}
  \item[RQ1] How do gender-related inventory responses vary across models,
        prompt structures and image conditions, and how do response coverage
        and general agreement affect their interpretation?
  \item[RQ2] How robust are models' image search judgements under scale
        reversal, and do their inventory scores correspond with those judgements
        relative to the historical human baseline?
\end{description}

RQ1 concerns responses to the adapted instrument, without assuming that they
establish an underlying psychological attitude. RQ2 examines descriptive
correspondence across five models from two families; this sample does not
support a general predictive relationship between inventory scores and search
judgements. The human comparison also retains the task and baseline differences
outlined above.

The four proposal aims were addressed to the following extent:

\begin{enumerate}
  \item The inventory audit was implemented for all five models, using prompt
        structures adapted from Gao and Kreiss \cite{gao2025measuring} and
        reading next-token probabilities before decoding. A face-photograph
        experiment extends the comparison to depicted gender
        (\cref{ch:study1}).
  \item The search evaluation uses the 2018 archive and the low-ASI human
        baseline. The contemporary recrawl and comparison with the high-ASI
        group were not included in the delivered analysis (\cref{ch:study2}).
  \item The relationship between inventory scores and downstream deviations
        is reported descriptively across five models, accounting for the
        small sample and dependence between family variants
        (\cref{sec:s2-deviation}).
  \item Scale reversal, adapted from CALM \cite{ye2025justice}, tests rating
        robustness with and without images. Image presence and referent
        wording change together, so the comparison does not isolate the
        visual pathway. CALM consistency rate, which concerns repeated
        judgements under identical conditions, is not separately reported
        (\cref{sec:framework-calm,sec:s2-calm}).
\end{enumerate}

% ---------------------------------------------------------------------
\section{Contributions}\label{sec:intro-contrib}
% \begin{plan}
% Five numbered contributions. Word them to match the \emph{final} results
% (several thesis-plan claims change after the new analyses --- see
% \cref{sec:s1-inference,sec:s2-deviation}):
% \begin{enumerate}
%   \item A logit-level audit protocol for VLMs that never reads generated
%         text (\cref{ch:framework}).
%   \item Captured mass $\capm$ as a validity statistic reported with every
%         measurement; with an output-diversity check it exposed silent
%         measurement faults that produced plausible aggregates
%         (\cref{tab:faults}).
%   \item Visual input collapses scale-reversal robustness in both production
%         models ($\dRR = -0.856$ and $-0.744$) (\cref{sec:s2-calm}).
%   \item The Llama models lose most of their yes/no mass whenever an image is
%         present, Qwen does not --- defensiveness is model-specific
%         (\cref{sec:s1-quality}).
%   \item $\ASIint$ is confounded by response style and the inventory's 16:6
%         keying imbalance; a shift-invariant alternative with an exact test
%         is given (\cref{sec:framework-responsebias}).
% \end{enumerate}
% \end{plan}

This thesis develops a logit-level audit protocol that reads next-token
probabilities in a single forward pass, without generating text. The protocol
specifies response-token pooling, chat templates and assistant prefills, and
reports captured mass alongside the resulting measurements
(\cref{ch:framework,ch:design}). Its application produces five main findings:

\begin{enumerate}
  \item Responses to face photographs differ by gender on the benevolent
        subscale. Signed scores are higher for female photographs in four
        models, with the opposite direction in Qwen-2B. Four of the five
        differences reach the corrected significance threshold, compared
        with only one combined-inventory difference. For Qwen-72B, conditional
        yes-probability on positively keyed benevolent items is $0.265$ for
        female photographs and $0.185$ for male photographs, excluding
        inversion prompts (\cref{sec:s1b-gap}).

  \item The signed aggregate inventory score is sensitive to general
        agreement. An exact decomposition shows that Llama-90B's score of
        $0.153$ is close to the $0.160$ predicted by its mean assent rate
        alone. Polarity separation complements this score by comparing
        positively and reverse-keyed items, remaining invariant to uniform
        shifts in agreement. All five models have negative separation, but
        none of the exact permutation tests survives Bonferroni correction
        across models (\cref{sec:framework-responsebias,sec:s1-inference}).

  \item Response coverage differs between the evaluated families when images
        are supplied. Llama-90B's captured mass falls from $0.951$ without
        an image to $0.467$ with a grey patch and $0.329$ with noise, while
        Qwen-72B remains near $0.95$. A featureless image is therefore
        sufficient to change Llama's coverage, without establishing refusal
        or algorithmic defensiveness. Coverage and output-diversity checks
        also helped expose silent measurement faults
        (\cref{sec:s1-quality,tab:faults}).

  \item The two largest models are less robust to scale reversal in the
        image condition. Robustness falls from $0.900$ without images to
        $0.044$ for Llama-90B and $0.156$ for Qwen-72B, giving
        $\dRR=-0.856$ and $-0.744$. The accompanying wording change prevents
        attributing the difference solely to visual input
        (\cref{sec:s2-calm}).

  \item The human comparison does not establish that the relationship
        between benevolent sexism and bias perception repeats across models.
        Three models rate ``smart person'' furthest below the low-ASI
        baseline, but the cross-model relationship remains inconclusive.
        The analysis distinguishes this query-level observation from
        replication of the human association and from a ranking of model
        attitudes (\cref{sec:s2-deviation,sec:s2-sigir}).
\end{enumerate}

% ---------------------------------------------------------------------
\section{Scope and thesis outline}\label{sec:intro-outline}
% \begin{plan}
% Half a page. One paragraph on scope changes from the proposal
% (contemporary 2026 stream not collected; attention heatmaps not
% produced; humanoid silhouette replaced by a grey patch) with a pointer to
% \cref{tab:scope}. Then one sentence per chapter.
% \end{plan}

The principal scope changes concern the historical image stream, attribution
analysis and Phase~1 visual anchor.

The contemporary 2026 image search stream was not collected. Phase~2 therefore
uses the 2018 archive alone, and the planned temporal comparison is absent.
Cross attention heatmaps were also not produced. The assessment of how visual
input affects the models consequently rests on behavioural comparisons
between conditions, without attribution to components within the model.
Finally, a uniform grey patch replaced the humanoid silhouette intended as
the Phase~1 visual anchor. The silhouette failed stimulus validation because
the 11B model returned a gender gap of $0.853$, above the threshold of $0.20$.
\cref{tab:scope} records the proposal and delivered scope, including
measurement counts and analyses not separately reported. The consequences for interpretation and generalisation
are discussed in \cref{sec:disc-threats}.

The remaining chapters develop the background, measurement framework and
experimental evaluation. \cref{ch:background} reviews bias in vision
encoders, the limitations of automated judges, response sensitivity following
alignment, and the human baselines motivating this thesis. It treats
algorithmic defensiveness as a hypothesis requiring investigation. It then
identifies the research gap. \cref{ch:framework} defines logit extraction,
captured mass, and the intrinsic and extrinsic metrics, together with the
propositions that establish their invariance properties.
\cref{ch:design} describes the models, instruments, stimulus construction
and validation, and the CALM design. It also documents the measurement
faults encountered during implementation.

\cref{ch:study1} reports the intrinsic study and answers RQ1, while
\cref{ch:study2} presents the downstream evaluation and addresses RQ2.
\cref{ch:discussion} brings the results together, explains the capabilities
and limitations of auditing at the level of logits, and tabulates the
threats to validity. \cref{ch:conclusion} summarises the contributions and
identifies further measurements needed to address the remaining questions.
